\chapter{Introduction}

\section{Background}
In the evolving landscape of digital healthcare, the precision of patient data is the bedrock of operational efficiency. This research explores the synthesis of biological scaling laws and computer vision to automate the pediatric census at Komfo Anokye Teaching Hospital (KATH).

\section{Problem Statement}
In major Ghanaian hospitals, optimizing medical staff allocation through the Local Health Information Management System (LHIMS) requires precise demographic data. However, the ``Invisible Child'' phenomenon—infants carried on mothers' backs—creates overlapping geometries that standard algorithms actively erase as visual noise (Addotey-Delove et al., 2023).

\section{Research Gap}
Specialized models like YOLOCDD have been developed for controlled medical scenes (Diop et al., 2025), but there is a need to extend these to high-density occlusion in African clinical norms.

\section{Research Objectives}
\begin{itemize}
    \item Formulate a spatial detection model to mathematically separate carried children from caregivers.
    \item Apply physical constraints via Head-to-Body Ratios to prevent misclassification.
    \item Implement geometric centroid tracking utilizing Kalman Filters for unique patient IDs.
    \item Design a localized Edge computing architecture that outputs strictly numerical matrices.
\end{itemize}
